\documentclass[11pt]{article}
\usepackage[a4paper,margin=23mm]{geometry}
\usepackage[T1]{fontenc}
\usepackage{lmodern,microtype,amsmath,amssymb,booktabs,tabularx,graphicx,hyperref,placeins}
\hypersetup{colorlinks=true,linkcolor=blue,citecolor=blue,urlcolor=blue}
\emergencystretch=2em
\newcommand{\pp}{\,pp}
\newcommand{\styunit}{mg\,MeOH\,h$^{-1}$\,g$_{cat}^{-1}$}
\newcommand{\plotfigure}[3]{\begin{figure}[htbp]\centering\includegraphics[width=\linewidth]{../figures/#1.pdf}\caption{#2}\label{#3}\end{figure}}
\title{Few-label calibration under publication shift in catalysis:\\what target labels contribute and where to acquire them}
\author{}
\date{}
\begin{document}\maketitle
\begin{abstract}
Few target-source measurements can recover catalytic prediction accuracy, but their value cannot be attributed to label count alone. We separate response level, input--label correspondence, historical response reuse and acquisition geometry under publication shift. Source-disjoint analyses cover methanol conversion, selectivity and space-time yield (STY), plus operating-only and composition-enriched oxidative coupling of methane (OCM): five panels from four outcomes in two reaction systems. A retrospective extension crosses five acquisition rules with thirteen predictors, retaining fixed pools, queries and nested label budgets. At five labels, established linear V-optimal acquisition reduces an unchanged blend's source-equal mean absolute error from 62.608 to 50.274 in native STY units and from 3.081 to 2.708 percentage points for composition-enriched OCM; both gains survive five-comparison adjustment. Conversion and selectivity show favorable point estimates but inconclusive adjusted intervals; operating-only OCM shows no gain. Previously evaluated additional OCM publications give a reduction from 3.552 to 3.140 percentage points. Label shuffling preserves response means yet degrades informative feature-dependent fits. A matched OCM follow-up separates acquisition representation from prediction representation. More elaborate historical covariance transfer does not improve strong controls. Labels support source-conditioned prediction, and informative locations can increase their value; neither establishes universal zero-label generalization or a superior new algorithm.
\end{abstract}

\section{Introduction}
Literature-derived catalytic measurements combine observations made with different materials, preparation procedures, reactors and reporting conventions. This heterogeneity is scientifically valuable: it gives access to operating ranges and catalyst families that a single experiment rarely covers. It also changes the prediction problem. A model that predicts another record from an already represented publication has received information about that source through related training records. A model deployed to an unseen publication has not. Both tasks are legitimate, but their validation results answer different questions.

Recent methanol collections provide catalyst and operating descriptors for CO$_2$ conversion, selectivity and productivity \cite{toldy,yan,suvarna}. OCM compilations connect elemental composition and reaction conditions with C$_2$ yield \cite{schmack,mine}. These resources make controlled cross-publication analysis possible. Grouped validation is established for dependent observations \cite{roberts}, and formula-grouped validation has also been examined for adsorption-energy screening \cite{wu}. Domain-adaptation theory motivates separating source-domain performance from target-domain reliability \cite{bendavid}; it does not supply a guarantee for the regressors or catalyst datasets studied here.

A practical response to publication shift is to acquire a small number of target-source measurements. Such labels can estimate a local response level, a residual offset or a within-source feature--response relationship. Calibration transfer has a long history in chemometrics \cite{feudale}. However, an improvement over an unadapted global model cannot alone demonstrate successful transfer of that model's response structure. A simple regressor fitted to the support responses may recover much of the same gain. Its comparison with a genuinely global-response-dependent adapter is therefore central.

We ask three related questions. How large is the prediction gap between record and source splits under controlled fitting? Which information accounts for the few-label gain: an output level, input--label correspondence, or additional global response structure? Do gains relative to a strong local response fit persist on additional publications from an independently released compilation? A fourth question asks whether label locations, at the same budget, change calibration value. The experiments separate label quantity, acquisition geometry, and predictor increments. A five-panel extension tests the portability of acquisition gains, with local, kernel and historical-response controls. More complex transfer candidates are evaluated as secondary mechanism checks.

The study contributes a source-disjoint comparison of label information and historical response reuse; projection and functional-response identities that expose support identifiability; and an equal-budget acquisition factorial that tests label geometry separately from predictor choice. An additional-publication experiment establishes a practical calibration gain while constraining stronger claims about response transfer. The estimators build on ridge regression, empirical-Bayes shrinkage and predictive combination \cite{efron,breiman}, with established functional decomposition and experimental-design ingredients. Their informative and unsuccessful combinations are retained to distinguish supported scientific findings from an unconfirmed new-algorithm claim.

\section{Data and evaluation design}
\subsection{Corpora, outcomes and source identity}
The historical analysis uses 1,245 paired methanol conversion/selectivity records from 82 publication sources and 80 leakage groups; 1,230 methanol space-time-yield (STY) records from 109 sources; and 1,794 OCM records from 408 publication entries and 405 leakage groups. The methanol conversion/selectivity corpus combines TheMeCat and the methanol retrospective \cite{toldy,yan}; STY uses the Suvarna public data \cite{suvarna}; OCM uses the Schmack supplement \cite{schmack}. These are two reaction systems and four outcomes. Operating-only OCM (F0) and composition-enriched OCM (F1) are paired representations of the same records, yielding five analysis panels. Twenty publication identifiers overlap the methanol conversion/selectivity and STY collections; those targets are not independent reaction replications.

A source is a publication entry, not a laboratory or physical catalyst batch. Historical duplicate citations and identified upstream overlaps are assigned to leakage groups. A group never crosses an outer source-disjoint split. OCM identities originate in numbered bibliographies, including some patents, so their provenance resolution is less uniform than DOI-normalized methanol sources. The released input snapshots preserve row and source identifiers. The supplementary information specifies missingness, preprocessing and the scope of the inherited corpus construction.

Conversion, selectivity and C$_2$ yield are scored in percentage points. STY is scored in \styunit. Methanol operating features are temperature, pressure, H$_2$/CO$_2$, GHSV and BET area. OCM operating features are temperature, CH$_4$ and O$_2$ partial pressures, their ratio, total pressure and contact time. F1 additionally encodes the reported cation and support species, known molar amounts and missing-amount indicators. Role-specific vocabularies are fitted on allowed training rows, with explicit channels for unseen species. Slots are order-invariant within a role; reported roles and amounts are not physically inferred or renormalized. These descriptors do not recover surface structure or preparation history.

\subsection{Controlled split audit}
To isolate split design, we refit two fixed historical-data configurations under record-level and leakage-group-disjoint five-fold splits, using seeds 42, 101, 202, 303 and 404. One is ridge regression with $\alpha=10$; the other is gradient boosting with 300 trees of depth 3. Neither is chosen by its audit test score. Imputation, scaling and composition vocabularies are refitted on the training partition in every case. The two model families are reported separately. This retrospective audit is distinct from the training-selected backbones used in the adaptation benchmark.

\subsection{Source-disjoint adaptation benchmark}
The adaptation study uses the same five outer seeds and five leakage-group folds. A global backbone is selected inside each outer-training partition by three-fold grouped out-of-fold (OOF) source-equal MAE. Its candidates are ridge $\alpha\in\{0.1,1,10,100\}$, gradient boosting with 100 or 300 trees and depth 2 or 3, and a 500-tree depth-8 candidate. Adapter settings are selected using historical training sources as pseudo-targets, with ten label-blind support draws per source and budget. All actual query outcomes remain outside selection. This includes the regularization choices for the local response baseline.

For each held-out source, query size is $n_T=\max(5,\lceil n_s/4\rceil)$, chosen by a stable identifier-based shuffle without labels. The remaining records form a support pool. The complete-curve cohort requires a pool of at least ten records: 32 sources for each paired methanol outcome, 22 for OCM and 28 for STY. Fifty shuffled support orders generate nested prefixes at $k\in\{1,2,3,5,10\}$. Methods share the query set and the same support prefix. The population is fixed across budgets; a lower MAE cannot be caused by dropping difficult sources at larger budgets.

Two secondary stress designs are also label-blind. High-temperature queries comprise an entire upper temperature block, with every support temperature strictly lower than every query temperature. A cutoff closest to a quarter of the records is selected among cutoffs leaving at least five queries and ten pool records. This yields 23 methanol, five historical OCM and 21 STY sources. The five-source historical OCM panel is exploratory. A reported-catalyst-name holdout is available for 16 methanol sources, where one complete name stratum is queried and other reported names supply support. It is a string-stratum test, not certification of novel physical catalysts. Hyperparameters are frozen from random-source tuning for both stress designs.

\subsection{Additional-source OCM test}
The external question and inclusion rules were locally frozen before new query outcomes were evaluated; this is not independent preregistration. We retrieved the Mine OCM update \cite{mine} at public repository commit \texttt{dbf7be5aadb30323f825e8f7d7f6bd3012c99dd2}. Its 4,759 records span 533 publication numbers. Comparison with the complete historical bibliography initially identified 121 additional numbers. Bibliographic screening then removed three re-listed old publications: new numbers 424, 427 and 429 correspond to old numbers 396, 395 and 105, respectively. This correction uses authors, journal, year, volume and pages, not prediction error.

We admit thermal OCM records with consistent units, attributable additional citations, valid measured yields and the required positive operating inputs. Missing contact time may use the historical training imputer. Records exceeding the frozen four-cation/two-support schema are excluded rather than silently discarding reported species. Full-record duplicates are removed and every excluded row retains its reason. The resulting external set contains 2,381 records from 91 additional publications; 52 sources support the fixed complete random budget curve and 40 support the high-temperature design. This establishes publication non-overlap in a separately released compilation, not independence of laboratories, extraction errors or physical batches, and not a strictly chronological test.

For the primary deployment test, preprocessing, global fitting and adapter tuning use all historical OCM records only, with the already specified training-selection procedure and seed 42. New sources supply support labels for adaptation but no tuning or query labels. The pre-specified primary contrast is F1 shrinkage blend versus local ridge at $k=5$. F0, other budgets, high-temperature results and subsequent mechanism diagnostics are secondary. External source filtering and all fitted parameters are archived. A post-hoc sensitivity additionally applies all 25 historical outer-training fits, without choosing among them on new outcomes, to examine dependence on the full-corpus refit regime.

\section{Adaptation methods and information content}
\subsection{Support-local and global residual responses}
Let $f_G$ be fitted on historical training sources, $A$ the labelled target-source support set and $T$ the disjoint query set. Write $p=f_G(x)$ and $z(x)$ for the training-fitted features. The support-local ridge prediction is
\begin{equation}
q_L(x)=b_L+z(x)^\top\beta_L,\quad
(b_L,\beta_L)=\arg\min_{b,\beta}\sum_{i\in A}(y_i-b-z_i^\top\beta)^2+\alpha_L\|\beta\|_2^2.
\end{equation}
Its response coefficients use only support responses; its preprocessing and regularization selection still use historical training information. It therefore removes direct use of a fitted global response predictor, not all forms of historical knowledge.

Residual ridge instead fits $y_i-p_i$ and retains the global prediction:
\begin{equation}
q_R(x)=f_G(x)+b_R+z(x)^\top\beta_R,\quad
(b_R,\beta_R)=\arg\min_{b,\beta}\sum_{i\in A}(y_i-p_i-b-z_i^\top\beta)^2+\alpha_R\|\beta\|_2^2.
\end{equation}
Intercepts are unpenalized. For fixed support/query features and the same $\alpha$, the ridge prediction operator $R_\alpha$ is linear in the response vector, including its fitted intercept. Hence
\begin{equation}\label{eq:identity}
q_R(T)-q_L(T)=f_G(T)-R_\alpha[f_G(A)](T).
\end{equation}
This identifies the global prediction structure not reconstructed by the support-conditioned linear projection. It is an information decomposition, not an error guarantee. Independently selected $\alpha_L$ and $\alpha_R$ also change regularization, so a matched-$\alpha$ ablation is required before assigning all observed gain to global structure. At $k=1$, $q_L$ equals the support label and $q_R$ equals the global prediction plus a mean residual offset.

\subsection{Shrinkage blending}
We combine the separately tuned components as $q_M=(1-w)q_L+wq_R$, $0\leq w\leq1$. For $k>1$, let $l_{-i},r_{-i}$ be support leave-one-out predictions and $d_i=r_{-i}-l_{-i}$. A training-source prior supplies $w_0$ and disagreement scale $D^2$. The support-only objective and solution are
\begin{align}
\widehat w&=\arg\min_{0\leq w\leq1}\sum_{i\in A}(y_i-l_{-i}-wd_i)^2+cD^2(w-w_0)^2,\\
\widehat w&=\operatorname{clip}_{[0,1]}
\frac{\sum_i d_i(y_i-l_{-i})+cD^2w_0}{\sum_i d_i^2+cD^2}.
\end{align}
For one label, infinite shrinkage or a zero denominator, the relevant prior weight is used. The prior grid is $w_0\in\{0,0.25,0.5,0.75,1\}$; $c\in\{\infty,10,2,0.5,0\}$. During pseudo-source shrinkage selection, the prior and disagreement scale exclude that pseudo-source's leakage group. The deployment prediction interface accepts support outcomes, but no query outcomes. Support leave-one-out evidence can stabilize a noisy small-sample mixture; it does not certify out-of-source reliability.

The other seven strategies are the unadapted global model, support mean, mean residual offset, penalized affine correction, empirical-Bayes residual offset, local ridge and temperature-only ridge. The full nine-way comparison distinguishes labels that estimate a level from labels that fit a feature-dependent response and from explicit global-response reuse. The affine and ridge grids use $\{0.1,1,10,100\}$. Empirical-Bayes hyperparameters are estimated from historical OOF residuals, rather than from held-out source queries. All formulas and exact leave-one-out implementation details are in the supplementary information.


\subsection{Functional responses and support identifiability}
The acquisition extension represents a frozen OCM F1 backbone through an empirical product-reference functional ANOVA. Let $O$ denote operating variables and $C$ the reported-composition encoding. A reference set contains one row from each of at most 32 historical training sources. With empirical marginal reference distributions, define
\begin{align}
c_0&=\mathbb E_{O_r,C_r} f_G(O_r,C_r),\\
u_O(x)&=\mathbb E_{C_r}f_G(O_x,C_r)-c_0,\quad
u_C(x)=\mathbb E_{O_r}f_G(O_r,C_x)-c_0,\\
u_I(x)&=f_G(x)-c_0-u_O(x)-u_C(x).
\end{align}
The three functions reconstruct the predictor exactly. Functional decomposition is established methodology \cite{hooker}; this product-reference implementation does not preserve the observed dependence of operating conditions and materials. Mixed reference inputs can involve unsupported combinations. These are predictor responses, not identified chemical causes.

For $U=(u_O,u_C,u_I)$, set $G_A(T)=U_T-H_{A,\alpha}U_A$, where $H_{A,\alpha}$ maps support responses to query ridge predictions, including an unpenalized intercept. The component predictor is $q=q_L+G_A\widehat w$, with box-constrained weights and a historical prior. Support leave-one-out evidence is $D_i=U_i-R_{A\setminus i,\alpha}[U_{A\setminus i}](x_i)$. The prior/penalty choices are inherited verbatim from the preceding component study; query outcomes do not choose them. Details of the direction-dependent penalty are in the supplementary information.

The leave-one-out residual operator annihilates constants, so $\operatorname{rank}(D)\leq\min(3,k-1)$. A composition-main response constant on the support set has a zero corresponding column in $D$. Those labels cannot independently identify its coefficient, even if the query composition differs. Predictions can still retain that response through a historical prior. This establishes a concrete limit of label information without promising that greater rank improves prediction.

\subsection{A label-blind innovation-risk acquisition candidate}
We investigate selecting label locations for the deployed correction operator, rather than maximizing information about a separately refitted Bayesian estimator. Let $V_j=u_j/s_j$, with historical response scales $s_j$. Under a working model $f_T=m_T+V_T\delta$ and $y_A=m_A+V_A\delta+\epsilon$, assume $\mathbb E\delta=0$, $\operatorname{Cov}(\delta)=B$, independent support noise $\operatorname{Cov}(\epsilon)=\sigma^2I$, and a fixed prior-mean correction $q_A=m_T+H_A(y_A-m_A)$. Its latent squared-error risk is
\begin{equation}\label{eq:designrisk}
\mathcal R(A)=\frac1{|T|}\operatorname{tr}\{(V_T-H_AV_A)B(V_T-H_AV_A)^\top\}
+\frac{\sigma^2}{|T|}\|H_A\|_F^2.
\end{equation}
The first term measures unresolved response directions on queries; the second penalizes support-noise amplification. The identity is exact under these assumptions for the fixed-prior predictor. For label-updated weights, it is a surrogate; it is not their exact risk. Neither the working-model assumptions nor covariance calibration follows from the identity.

We estimate $B$ and $\sigma^2$ exclusively from historical source-specific regressions on cross-fitted component responses. The covariance combines half the across-source covariance, half its diagonal, and mean within-regression posterior uncertainty. Residual variance is a lack-of-fit/noise proxy, not independently measured experimental noise. The selector greedily minimizes Eq.~\ref{eq:designrisk} for each nested acquisition prefix, using only pool/query covariates and frozen historical quantities. No support or query outcomes are inspected to choose positions. No global optimality, monotonic-risk, or real-data generalization guarantee is claimed.

\subsection{Crossing acquisition and prediction}
The retrospective OCM F1 extension crosses ten acquisition policies with thirteen predictors. Controls include random acquisition, temperature/feature maximin coverage, and ordinary greedy integrated posterior-variance (V-optimal) selection with linear and component-augmented features. Transductive experimental design and variance-based acquisition are established approaches \cite{yubi,fontaine}; their principles are not claimed as new. Four innovation-risk ablations replace the response covariance by an isotropic one, remove component uncertainty, remove the noise term, or target the pool rather than known query covariates. A nuisance-parameter distinction is also established in Bayesian design \cite{sloman}; our three-response correction construction is a specific working model to be evaluated.

All methods receive the identical support prefix within a policy and budget; different policies share the same fixed query set and eligible pool. All 25 historical outer fits and the additional-publication fit are reused under their original training exclusions. Frozen local ridge, original blend, RBF and ordinary two-component feature-augmented ridge are retained. Additional controls retune local regularization, RBF settings, blend prior weight, and a three-component ordinary augmented ridge on historical pseudo-sources for each policy/budget. The blend's two ridge regularizers and shrinkage strength remain inherited; this is a partial, stated retuning, not exhaustive optimization of every estimator.

Random acquisition uses fifty support orders, while each deterministic policy uses one nested order. Repetition counts are averaged within each source and policy before seed/source averaging. The ten-by-thirteen combinations are not 130 statistically independent algorithms. Query-informed policies assume query covariates are available before acquisition, which defines a transductive deployment setting. The corpus has already been inspected in earlier work, so the acquisition study is retrospective even though target query outcomes remain excluded from tuning and selection. Most intervals are nominal exploratory source-group intervals. Four candidate-focused contrasts at random-query $k=5$ additionally receive Bonferroni-adjusted percentile intervals; this adjustment does not convert the development corpus into an untouched test.

\subsection{Metrics and uncertainty}
Within each source, we average query absolute error over support draws and then outer seeds; sources receive equal final weight. External full-corpus deployment has one frozen historical fit, so only support draws are averaged before sources. Record-weighted MAE and physically bounded predictions are secondary sensitivities. Paired method differences use the same source, support and query population; negative differences favor the left-hand method. We resample leakage groups 20,000 times, retaining source weights within groups. Intervals are conditional on the fitted models and observed source population, with no simultaneous correction across the many secondary comparisons. Neither support draws nor folds are counted as independent publications.

\subsection{Secondary uncertainty-transfer candidates}
Component-aware acquisition and a subsequent residual-covariance candidate test whether historically learned uncertainty adds predictive value. The latter is a finite-rank Gaussian-process/random-coefficient model with a free source intercept, historically learned coefficient covariance, and aligned posterior-variance acquisition. These ingredients are established \cite{gpml,alpaca,gpmllearn,scaml}. Source-disjoint response features and a further nested cross-fit exclude the entire pseudo-validation fold during covariance learning. Five covariance alternatives and strong policy-specific controls are evaluated, with four locked candidate contrasts per cohort. Complete derivations, parameter-matched ablations and all outcomes are retained in Supplementary Section 13. Working covariance is not measured laboratory noise or validated predictive uncertainty \cite{iwata}.

\subsection{Five-panel test of acquisition geometry}
The new retrospective extension uses every one of the 25 inherited outer fits per panel, with the original query populations and budgets $1,2,3,5,10$. Nine frozen predictors are crossed with random acquisition, temperature coverage, full-feature coverage, query-directed linear V-optimal design and a pool-only counterpart. Four additional controls re-tune local and residual ridge, blending, and local RBF regression by policy and budget on historical pseudo-sources only. Historical source-disjoint OOF predictions are reused with their matching inner training preprocessors; 125 outer frozen-configuration models are actually re-fitted. This extension does not train 125 new nested searches or choose configurations from outer queries.

For a working feature row $\psi(x)=[10,z(x)^\top]^\top$, the standard linear design objective is
\begin{equation}
J_Q(A)=\frac{1}{n_Q}\operatorname{tr}\left[F_Q\left(I-F_A^\top(F_AF_A^\top+\alpha I)^{-1}F_A\right)F_Q^\top\right].
\label{eq:crossdesign}
\end{equation}
Each greedy addition minimizes $J_Q$. Its unit coefficient prior, intercept prior variance 100 and inherited ridge regularizer form a working surrogate, not a claim of calibrated physical uncertainty. Prediction retains the original unpenalized intercept. Query V uses known query covariates and is transductive; pool V replaces $F_Q$ with the candidate pool and does not use query covariates. Coverage rules choose the point nearest the pool center first, then maximize distance from selected points. Every selector is label-blind. Random uses 50 nested orders; each deterministic rule uses one. All predictors share a policy's supports. V-optimal design is established \cite{yubi,fontaine}; these are explicit standard-design implementations, not complete reproductions of those papers' algorithms.

The primary family is the frozen blend's query-V-minus-random MAE at five labels and random-query allocation in each of the five panels. Leakage-group bootstrap uses 20,000 draws, averaging support draws and seeds before sources. Nominal 95\% and five-member Bonferroni-adjusted 99\% percentile intervals are reported. The panels share records and sources, so they are not independent replications; Bonferroni adjustment does not require independence. All other budgets, rules, control increments, stress panels and label shuffles are exploratory. This local protocol lock precedes current outer evaluation, but the same corpora have already informed earlier method development.

\section{Results}
\subsection{Record-level accuracy does not describe publication transfer}
The fixed gradient-boosting audit shows lower stacked $R^2$ under source splits in every panel (Fig.~\ref{fig:split}). Conversion changes from 0.794 to 0.236; methanol selectivity from 0.679 to $-0.002$; STY from 0.782 to 0.344; OCM F0 from 0.314 to 0.001; and OCM F1 from 0.458 to 0.235. The observed drops range from 0.223 to 0.681. Fixed ridge is less accurate for several outcomes and its split gaps are smaller, showing why optimism should be reported with the fitted model family rather than treated as a dataset constant. Full source- and record-weighted errors are supplied for both families.
\plotfigure{fig1_split_audit}{Controlled fixed-configuration audit. Each line contrasts record-level and source-disjoint five-fold prediction under the same model configuration and training-only transformations. $R^2$ pools stacked repeated out-of-fold predictions; it is not $R^2$ of the averaged prediction. F0 and F1 share OCM rows.}{fig:split}

\subsection{Internally, labels recover accuracy but extra transfer is task-dependent}
Figure~\ref{fig:budget} and Table~\ref{tab:internal} report all strategies on the fixed historical cohorts. At five labels, local ridge gives 2.429\pp\ conversion, 11.547\pp\ selectivity, 3.655\pp\ OCM F0, 3.374\pp\ OCM F1 and 65.210 STY units. The composition-enriched OCM blend gives 3.081\pp, a paired difference of $-0.294$ [$-0.449,-0.154$]\pp\ from local ridge, improving 18 of 22 sources. Its average benefit is present at all five budgets, with paired intervals below zero; these are retrospective internal findings, not five independent confirmations.

Residual ridge alone improves F1 by $-0.265$ [$-0.485,-0.056$]\pp\ at five labels. In contrast, it worsens conversion relative to local ridge by 0.291 [0.075,0.500]\pp\ and does not show a clear advantage for selectivity or STY. The blend remains close to local ridge in the paired methanol outcomes and OCM F0. This pattern contradicts a universal preference for retaining a historical response model.
\begin{table}[htbp]\centering\small
\caption{Historical random-source cohorts at $k=5$, source-equal MAE. Percent outcomes use pp; STY uses \styunit. Global predictions use no support labels.}\label{tab:internal}
\begin{tabular}{lrrrrr}
\toprule
Strategy & Conv. & Sel. & OCM F0 & OCM F1 & STY \\
\midrule
Global & 5.271 & 22.740 & 5.302 & 4.527 & 139.258 \\
Support mean & 4.304 & 17.136 & 3.693 & 3.693 & 91.993 \\
Offset & 3.455 & 15.747 & 3.672 & 3.423 & 88.265 \\
Affine & 3.083 & 12.539 & 3.694 & 3.360 & 71.225 \\
EB offset & 3.457 & 15.988 & 3.759 & 3.415 & 86.510 \\
Local ridge & 2.429 & 11.547 & 3.655 & 3.374 & 65.210 \\
Temperature ridge & 2.519 & 11.535 & 3.677 & 3.677 & 87.098 \\
Residual ridge & 2.720 & 12.077 & 3.701 & 3.109 & 69.889 \\
Shrinkage blend & 2.441 & 11.584 & 3.660 & 3.081 & 62.608 \\
\bottomrule
\end{tabular}
\end{table}
\plotfigure{fig2_internal_budgets}{Nine strategies under common query sets and nested support prefixes. Cohorts comprise 32 paired methanol sources, 22 OCM sources and 28 STY sources. Lines summarize source-equal errors; uncertainty and full per-source values are provided in the supplementary tables and machine-readable results.}{fig:budget}

\subsection{Additional publications confirm calibration value, not blend superiority}
On 52 additional OCM sources, composition-enriched global MAE is 6.120\pp. Five-label local ridge gives 3.593\pp, reducing error by 41.3\%; the paired difference is $-2.528$ [$-3.328,-1.763$]\pp\ and 42 sources improve. The support mean gives 4.490\pp, while temperature-only ridge gives 3.779\pp. Label information therefore supports more than a mean shift, and a multi-feature local response fit improves on a strong temperature baseline.

The frozen primary blend gives 3.552\pp. Its difference from local ridge is $-0.041$ [$-0.140,0.045$]\pp, with 21 improved and 31 harmed sources. A slightly better mean alongside more harmed sources reflects heterogeneous effect magnitudes. The interval is compatible with a modest benefit and a modest harm; it does not establish equivalence or reproduce the internal advantage. Across budgets, blend point improvements remain small and every interval includes zero (Fig.~\ref{fig:external}). The operating-only residual adapter is worse than local ridge at $k=5$: $+0.376$ [0.148,0.671]\pp. The corresponding blend stays close to local prediction.
\begin{table}[htbp]\centering\small
\caption{Additional-source OCM at $k=5$, source-equal MAE in pp. Random panels use 52 sources; high-temperature panels use 40. Inputs, support labels and query sets are paired between F0 and F1.}\label{tab:external}
\begin{tabular}{lrrrr}
\toprule
Strategy & F0 random & F1 random & F0 high T & F1 high T \\
\midrule
Global & 5.973 & 6.120 & 5.699 & 5.140 \\
Support mean & 4.490 & 4.490 & 5.556 & 5.556 \\
Offset & 4.502 & 4.606 & 5.512 & 5.397 \\
Affine & 4.284 & 4.410 & 5.209 & 5.328 \\
EB offset & 4.545 & 4.621 & 5.205 & 4.942 \\
Local ridge & 4.045 & 3.593 & 4.801 & 4.365 \\
Temperature ridge & 3.779 & 3.779 & 4.697 & 4.697 \\
Residual ridge & 4.421 & 3.716 & 5.375 & 4.454 \\
Shrinkage blend & 4.054 & 3.552 & 4.806 & 4.335 \\
\bottomrule
\end{tabular}
\end{table}
\plotfigure{fig3_external_budgets}{Frozen deployment on additional OCM publications. All historical fitting precedes support/query evaluation, and new query labels are used only for scoring. The small differences between blend and local ridge do not confirm the larger internal increment.}{fig:external}

\subsection{Representation and label correspondence explain distinct gains}
Composition reduces local OCM MAE relative to F0 by $-0.281$ [$-0.518,-0.056$]\pp\ internally and $-0.453$ [$-0.710,-0.229$]\pp\ on additional sources (Fig.~\ref{fig:representation}). Yet adding composition to the unadapted global model changes external MAE by $+0.148$ [$-0.383,0.675$]\pp. Richer inputs can help a calibrated local response without establishing stronger unadapted publication transfer. Support means and temperature-only predictions are invariant to this paired representation change.

To distinguish a response level from correspondence, we shuffle labels within the same five-record support set, retaining its empirical mean and feature set. The exploratory external diagnostic uses ten support draws per source and twenty permutations per draw, without retuning. For F1 local ridge, shuffled-label error exceeds real-label error by 1.265 [0.877,1.694]\pp; the F0 penalty is 0.502 [0.286,0.756]\pp. Blend and residual F1 predictions show similar correspondence penalties (Fig.~\ref{fig:alignment}). Mean, offset and empirical-Bayes offset controls are unchanged to numerical precision. These diagnostics support the importance of input--response pairing in calibration. They are not causal attribution or an automatically valid permutation significance test, and their ten-draw real arm is distinct from the fifty-draw headline estimate.
\plotfigure{fig5_label_alignment}{External label-alignment diagnostic. Positive values are the increase in MAE after support labels are permuted. Permutation preserves the support mean. Group-bootstrap intervals condition on fixed fitted models. Mean-only controls have approximately zero change and are reported in the supplementary table.}{fig:alignment}
\plotfigure{fig6_representation}{Paired F1 minus F0 MAE at five labels. Negative values favor the composition-enriched representation. Internal and external contrasts use identical records within their own cohort; they are not paired across cohorts. Representation benefits for local responses and for unadapted global models need not coincide.}{fig:representation}

\subsection{Transfer boundaries persist under stress and fit-regime checks}
For historical STY high-temperature queries, blending reduces five-label MAE from 133.260 to 115.145, a 13.6\% decrease. Its paired difference is $-18.114$ [$-39.291,-3.159$] STY units, but only 11 of 21 sources improve, again indicating heterogeneity. The random-query STY difference is $-2.602$ [$-6.193,0.323$], which is inconclusive. Reported-name holdout does not show a clear additional blend benefit for methanol. The supplementary results retain all stress panels, including the small historical OCM panel.

For the 40 additional OCM high-temperature sources, F1 local and blend MAEs are 4.365 and 4.335\pp; the paired difference is $-0.031$ [$-0.164,0.097$]\pp. Using all 25 historical outer-training fits instead of the full-corpus refit also fails to establish blend superiority: on random queries the difference is $-0.062$ [$-0.160,0.024$]\pp. This sensitivity does not prove a cause for the boundary, but makes a change of training-fit regime alone an insufficient explanation of the observed result.

Matched regularization further narrows interpretation. Setting residual $\alpha$ to the already selected local $\alpha$ gives an internal F1 difference of $-0.204$ [$-0.425,0.004$]\pp\ and an external difference of $+0.123$ [$-0.071,0.316$]\pp\ (Fig.~\ref{fig:matched}). The identity in Eq.~\ref{eq:identity} holds numerically on every evaluated support/query draw. Thus the separately tuned internal residual result cannot all be assigned to a statistically established global-structure effect. On the additional-source cohort, only three sources contain species unseen by the fitted role vocabulary and nine include operating values outside at least one historical range. Simple species novelty or range checks are therefore an incomplete description of publication shift; the present data do not identify its physical cause.
\plotfigure{fig4_transfer_contrasts}{Five-label residual and blend differences from local ridge. OCM and STY are displayed in separate panels with their own units. The number in parentheses is the source count. Average improvement can coexist with harm in many individual sources.}{fig:contrasts}
\plotfigure{fig7_matched_alpha}{Exploratory same-regularization diagnostic at $k=5$. Residual ridge uses the local method's frozen selected $\alpha$; there is no new tuning on queries. Intervals include zero in every panel. This ablation isolates the prediction component in Eq.~\ref{eq:identity}, rather than asserting guaranteed benefit from it.}{fig:matched}
\FloatBarrier


\subsection{Better label locations improve a fixed existing predictor}
The acquisition factorial separates where labels are taken from how their responses are fitted. With the original blend and all its parameters unchanged, linear V-optimal acquisition reduces five-label MAE from 3.081 to 2.708\pp\ internally and from 3.552 to 3.140\pp\ on the 52 additional sources (Table~\ref{tab:acquisition}). The matched source differences are $-0.372$ [$-0.650,-0.122$] and $-0.412$ [$-0.610,-0.211$]\pp, corresponding to reductions of 12.1\% and 11.6\%. Sixteen of 22 historical and 42 of 52 additional sources improve. This is a sampling increment with the same five-label budget; it is not attributed to a new transfer predictor.

Augmented V-optimal acquisition gives 2.737 and 3.153\pp\ for the same frozen blend, with matched differences of $-0.343$ [$-0.634,-0.074$] and $-0.398$ [$-0.602,-0.197$]\pp. Retuning the blend prior by policy retains the sampling gains: linear-design MAEs are 2.716 and 3.145\pp, versus random-design values of 3.084 and 3.552\pp. The available evidence supports informed acquisition in this retrospective population. It does not identify linear rather than augmented design as a universally optimal rule.
\begin{table}[htbp]\centering\small
\caption{Same five-label budget and frozen original blend under different acquisition rules. MAE uses source weights and pp. Sampling effects, rather than a predictor change, explain these comparisons.}\label{tab:acquisition}
\begin{tabular}{lrr}\toprule Acquisition & Historical ($n=22$) & Additional ($n=52$)\\\midrule
Random & 3.081 & 3.552\\
Temperature cover & 3.407 & 4.468\\
Feature cover & 2.921 & 3.405\\
Linear V-optimal & 2.708 & 3.140\\
Augmented V-optimal & 2.737 & 3.153\\
Innovation risk & 3.192 & 3.662\\
\bottomrule\end{tabular}\end{table}
\plotfigure{step5_02_sampling_contrasts}{Acquisition increments for the same frozen blend. Negative values favor the indicated rule over random support acquisition. Intervals are nominal source-group bootstrap intervals; the retrospective comparisons are conditional on fitted models.}{fig:acquisition}

\subsection{Component-aware risk does not outperform standard design}
The innovation-risk candidate gives ANOVA-transfer MAEs of 3.237\pp\ internally and 3.716\pp\ on additional random-query sources. Relative to random acquisition for that predictor, its changes are $+0.072$ [$-0.110,0.271$] and $+0.163$ [$-0.135,0.505$]\pp. Relative to augmented V-optimal acquisition with the same predictor, the candidate is worse by $+0.441$ [0.146,0.768] and $+0.557$ [0.252,0.888]\pp. The adjusted four-contrast intervals for the latter differences remain above zero. Thus neither an exact working-model risk identity nor a response decomposition establishes acquisition superiority.

At additional-source $k=5$, the inherited component penalty is infinite: ANOVA-transfer predictions equal fixed-prior predictions under every acquisition rule. Five selected labels therefore do not demonstrate learned component reweighting in this panel. On the linear-design support set, the internal ANOVA predictor is worse than the original blend by 0.101 [0.030,0.182]\pp; the additional-source difference is 0.001 [$-0.068,0.059$]\pp. Better supports help several predictors without validating the new component predictor's incremental originality.

A secondary squared-error diagnostic also exposes working-risk miscalibration. Internally, innovation-risk acquisition lowers the mean working risk to 11.117\pp$^2$ versus 14.209 for linear V-optimal, yet observed fixed-prior MSE is 20.144 versus 14.158\pp$^2$. The corresponding additional-source values are working risks 10.071 versus 10.612, and observed MSEs 32.610 versus 21.073\pp$^2$. The discrepancy cannot be explained solely by optimizing squared rather than absolute error. The data do not separately identify covariance error, residual/noise confounding, and function misspecification; each remains a plausible explanation.

\subsection{Acquisition benefits have a defined stress-test scope}
For 40 additional high-temperature sources, linear V-optimal acquisition reduces the frozen blend's five-label MAE from 4.335 to 3.563\pp. The matched change is $-0.772$ [$-1.205,-0.332$]\pp, or 17.8\%, with 34 sources improving. Augmented design gives a reduction of 15.0\%. The five-source historical high-temperature panel does not show a clear benefit: linear-design change is $+0.102$ [$-0.379,0.546$]\pp. These outcomes support a useful, explicitly bounded acquisition result, rather than a claim of uniform improvement across all shifts or budgets. Full factorial tables and budget curves retain the unsuccessful candidates and stress results.
\plotfigure{step5_03_budgets}{Frozen-blend budget curves under shared query populations and nested acquisition prefixes. The historical high-temperature panel has only five sources. The factorial extension is retrospective; standard design gains and the unsuccessful risk candidate remain visible together.}{fig:acquisitionbudget}

\subsection{Published-workflow audits and corrected statistical weighting}
We retain a controlled audit of Suvarna's public workflow and refit the Xu fixed-configuration grouping contrast \cite{suvarna,xu}. Their purpose is to examine evaluation scope, rather than to judge the credibility of the original authors or validate the new adapters. Re-aggregation of archived Suvarna XGBoost row predictions gives source-equal MAE 57.897 under record splits and 124.043 under source splits; stacked raw-STY $R^2$ changes from 0.876 to 0.319. This release independently verifies those archived errors without claiming an original-model bitwise refit.

The Xu main audit uses 1,757 DOI-attributable rows from 156 sources. We refit the same fixed XGBoost configuration under five record and source folds across five seeds in one current environment. Correct source weighting first pools all held-out records within each source and seed, then averages seeds and sources. The resulting MAEs are 69.620 and 126.365, respectively (ratio 1.815). Directly averaging source--fold cells instead gives 63.271 on the record side because sources appear in unequal numbers of folds. That superseded quantity is not a source-equal MAE. An archived random-80/20 reporting value is separately corrected to 62.978; it is not mixed with the five-fold estimate. These corrections preserve the source-transfer conclusion while making the statistical estimand explicit.

\subsection{Standard acquisition gains extend across tasks, with limits}
At the same five-label budget and with the same frozen blend, query V-optimal acquisition reduces STY MAE by 19.7\% and composition-enriched OCM MAE by 12.1\% (Table~\ref{tab:crossprimary}; Fig.~\ref{fig:crossprimary}). Their adjusted differences are $-12.335$ [$-25.965,-2.874$] in native STY units and $-0.372$ [$-0.746,-0.056$]\pp, respectively. The conversion point estimate improves 13.1\%, but its adjusted interval crosses zero; the nominal interval is favorable. Selectivity improves 7.8\% in point estimate with both nominal and adjusted intervals crossing zero. Operating-only OCM has a 3.0\% higher point estimate, also inconclusive. Four favorable points out of five therefore do not establish universal acquisition superiority.

\begin{table}[htbp]\centering\small
\caption{Locked five-panel primary comparison at five labels: unchanged frozen blend, random-query allocation. Gain is relative MAE reduction; negative values mean higher error. MAE and difference intervals use percentage points except STY, which uses \styunit. Intervals are adjusted across the five comparisons.}\label{tab:crossprimary}
\begin{tabular}{lrrrrl}\toprule Panel & $n$ & Random & Query V & Gain (\%) & Adjusted $\Delta$ interval\\\midrule
Conversion & 32 & 2.441 & 2.121 & 13.1 & [-0.707, 0.068]\\
Selectivity & 32 & 11.584 & 10.683 & 7.8 & [-2.140, 0.615]\\
STY & 28 & 62.608 & 50.274 & 19.7 & [-25.965, -2.874]\\
OCM F0 & 22 & 3.660 & 3.768 & -3.0 & [-0.355, 0.714]\\
OCM F1 & 22 & 3.081 & 2.708 & 12.1 & [-0.746, -0.056]\\
\bottomrule\end{tabular}
\end{table}
\plotfigure{step7_01_primary}{Standard query-directed V-optimal acquisition minus random acquisition. Intervals are the primary adjusted native-unit difference intervals rescaled by each panel's fixed observed random MAE for display; they are not a bootstrap ratio interval or a pooled reaction effect. Negative values favor query V. Source counts and representation dependence constrain the claim.}{fig:crossprimary}

The broader controls preserve this interpretation. At five labels, re-tuned blending gives STY MAE 49.852 under query V versus 62.613 under random acquisition. For OCM F1 it gives 2.696 versus 3.073\pp. On the same query-V supports, ordinary local ridge and local RBF give STY MAEs 51.293 and 51.547, and OCM F1 values 2.823 and 2.851\pp. Most blend increments over these local controls remain inconclusive in the other panels. Full-feature coverage and pool-only design can help, but neither rule is a universal winner. Pool-only selectivity has a lower point estimate than query V, and their paired interval crosses zero. The full factorial and budget curves are in Supplementary Section 14.

High-temperature panels also differ. Historical conversion and selectivity show exploratory query-V improvements on 23 sources, whereas STY's 21-source interval crosses zero. The historical OCM high-temperature set contains only five sources and does not show a reliable gain. These stress outcomes share query populations across acquisition rules; they are not pooled with random-query data or additional-corpus results.

\subsection{Labels encode correspondence, and representation also guides selection}
Under query V, five-label correspondence shuffling leaves the same label multiset and mean but changes the frozen blend's STY MAE from 50.274 to 80.945, and OCM F1 MAE from 2.708 to 3.548\pp. Conversion and selectivity also worsen. For OCM F0, the change is only $+0.047$\pp\ with a nominal interval crossing zero. These exploratory controls show where feature--response correspondence contributes beyond an output mean; they do not identify a causal chemical mechanism.

After observing the primary outcomes, we added a clearly post-primary OCM diagnostic crossing acquisition geometry and prediction representation. Operating-only and composition-enriched V design use the same inherited F1 regularizer schedule and the same pools/queries; F0 and F1 frozen predictors each see both support sets. With the F1 predictor, operating geometry gives 3.300\pp\ and composition geometry 2.708\pp; the nominal paired change is $-0.592$ [$-1.347,-0.020$]\pp. With the F0 predictor, these values are 3.821 and 3.370\pp, a nominal change of $-0.452$ [$-0.982,-0.010$]\pp. Composition features can thus carry useful selection information even for a predictor that does not consume those features. This is exploratory evidence from 22 sources, not an adjusted primary result or proof that composition is necessary in every reaction. STY demonstrates a useful operating-feature acquisition case.

\subsection{Complex uncertainty transfer remains a boundary result}
The preceding residual-covariance candidate does not provide a stronger algorithm. Its historical/additional five-label MAEs are 3.287/3.903\pp\ versus 2.696/3.149\pp\ for re-tuned blending with standard linear design. The four-comparison adjusted candidate increments are positive in both cohorts. Cross-fitted covariance improves a small matched in-sample-covariance contrast, but does not reverse overall inferiority. Working response variance also differs from observed error. Complete candidate results and unfavorable alternatives remain in Supplementary Section 13; their coherent formulas are not treated as performance innovation.

\section{Discussion}
\subsection{What does it mean for labels to improve generalization?}
Few-label calibration changes the task from unadapted transfer to prediction conditional on target-source measurements. The new labels can identify a response level and, when the features vary usefully within that source, a local response trend. They do not establish improved zero-label performance on arbitrary future sources. This distinction matters both for deployment and for interpreting a method's contribution.

The external experiment provides evidence that labels are useful: a fitted local response improves on the global predictor and the support mean, while shuffling labels makes feature-dependent calibration worse. It also provides evidence about their limits: additional global-response reuse contributes much less than the original internal estimate suggested. Local fitting and global transfer must therefore be evaluated as separate increments. A large recovery relative to the unadapted model should not be presented as if the same amount were attributable to the global model's transferable structure.

\subsection{Representation is necessary evidence, not a transfer certificate}
Without adequate catalyst descriptors, identical or nearly identical operating inputs may correspond to distinct material responses. Composition can reduce this ambiguity and help support labels choose a locally relevant relationship. However, composition categories do not encode surface state, synthesis, reactor geometry or measurement practice. The external result demonstrates that a representation useful for local calibration need not reliably improve an unadapted historical predictor. Equation~\ref{eq:identity} explains where extra global response information enters; only held-out error establishes whether that information helps.

\subsection{Practical implications and limitations}
The benchmark supports a source-local response model as a necessary deployment baseline. A residual or blend adapter can be worthwhile, but its incremental value should be measured on source-disjoint and separately compiled sources rather than inferred from a recovery curve against an unadapted model. We retain the original frozen primary result and separately identify the acquisition extension as retrospective development. Standard design improves the frozen blend on STY and composition-enriched OCM, with inconclusive adjusted conversion/selectivity results and no operating-only OCM gain. Both component-aware risk and a subsequently evaluated source-contrast covariance candidate fail to add a reliable benefit; the latter is worse than standard linear-design blending in the locked comparisons. The blend is not a validated universal recommendation or a safety guarantee; its prior can still carry an unreliable global component.

Several boundaries remain. Publications are imperfect domains; neither the historical nor additional-source data certify laboratory or batch separation. Eligibility favors sources with at least fifteen compatible records and the selected input schema. Historical catalyst names are incomplete, and STY inputs here are operating descriptors, although richer public descriptors exist. The external corpus broadens OCM publications within the same reaction, rather than adding an independent third reaction. The secondary analyses are retrospective and their intervals are conditional; the earlier four-contrast candidate families are adjusted separately within each cohort, and the current five-panel acquisition primary family also receives adjustment. The study evaluates prediction and offline selection of labels from a finite measured corpus, rather than causal chemical mechanisms or a prospective wet-laboratory campaign. A prospective campaign with catalyst/batch identifiers and a locked acquisition rule would test a stronger claim about experimental use.

\section{Conclusion}
Target-source labels supply an output level and, in informative representations, input--response correspondence. Acquisition geometry changes how much those labels contribute at the same budget. Standard query-directed design improves STY and composition-enriched OCM under a fixed predictor, while other panels set clear limits. An exploratory representation crossing also locates information useful for selection separately from information consumed by prediction. Historical response transfer contributes an additional, context-dependent increment; complex covariance repairs do not establish superiority. These findings support a bounded scientific claim about target information and measurement placement, not universal zero-label generalization or a new optimal-design algorithm.

\section*{Data and code availability}
The seventh numbered package contains cleaned multi-task input snapshots, all 125 inherited parameter/OOF/split records needed by this extension, current support/query ledgers, every current float64 query prediction, measured summary tables, code and complete manuscript/figure sources. Its entry point verifies released data and actually re-fits current outer models in a fresh directory; inherited OOF responses are explicitly reused for historical tuning. Earlier complete releases retain upstream attributions, additional-source inclusion/bibliography ledgers, external prediction exports and the full older candidate stages. Their identifiers and cryptographic hashes are recorded in the current package. A final public repository or archive DOI should integrate those stages; a local stage lock is not independent preregistration.

\begin{thebibliography}{12}\small
\bibitem{toldy} Toldy, \'A. I. et al. TheMeCat: Dataset of Thermocatalytic Conversion of CO$_2$ to Methanol. \emph{Scientific Data} \textbf{13}, 165 (2026). \url{https://doi.org/10.1038/s41597-025-06482-8}.
\bibitem{yan} Yan, M. et al. A machine learning perspective on three decades of methanol synthesis: research framework and experimental operation insights. \emph{Communications Engineering} \textbf{5}, 162 (2026). \url{https://doi.org/10.1038/s44172-026-00706-4}.
\bibitem{suvarna} Suvarna, M., Ara\'ujo, T. P. \& P\'erez-Ram\'irez, J. A generalized machine learning framework to predict the space-time yield of methanol from thermocatalytic CO$_2$ hydrogenation. \emph{Applied Catalysis B: Environmental} \textbf{315}, 121530 (2022). \url{https://doi.org/10.1016/j.apcatb.2022.121530}.
\bibitem{schmack} Schmack, R. et al. A meta-analysis of catalytic literature data reveals property--performance correlations for the OCM reaction. \emph{Nature Communications} \textbf{10}, 441 (2019). \url{https://doi.org/10.1038/s41467-019-08325-8}.
\bibitem{mine} Mine, S. et al. Analysis of Updated Literature Data up to 2019 on the Oxidative Coupling of Methane Using an Extrapolative Machine-Learning Method to Identify Novel Catalysts. \emph{ChemCatChem} (2021). \url{https://doi.org/10.1002/cctc.202100495}. Data: \url{https://github.com/mts-uw/OCM}, commit \texttt{dbf7be5aadb30323f825e8f7d7f6bd3012c99dd2}.
\bibitem{roberts} Roberts, D. R. et al. Cross-validation strategies for data with temporal, spatial, hierarchical, or phylogenetic structure. \emph{Ecography} \textbf{40}, 913--929 (2017). \url{https://doi.org/10.1111/ecog.02881}.
\bibitem{wu} Wu, W., Yang, M., Cheng, P. \& Wang, Y. Reliable Machine Learning Screening of Adsorption Energies Is Better Assessed with Formula-Grouped Cross-Validation. \emph{Catalysts} \textbf{16}, 742 (2026). \url{https://doi.org/10.3390/catal16080742}.
\bibitem{bendavid} Ben-David, S. et al. A theory of learning from different domains. \emph{Machine Learning} \textbf{79}, 151--175 (2010). \url{https://doi.org/10.1007/s10994-009-5152-4}.
\bibitem{feudale} Feudale, R. N. et al. Transfer of multivariate calibration models: a review. \emph{Chemometrics and Intelligent Laboratory Systems} \textbf{64}, 181--192 (2002). \url{https://doi.org/10.1016/S0169-7439(02)00085-0}.
\bibitem{efron} Efron, B. \& Morris, C. Stein's estimation rule and its competitors---an empirical Bayes approach. \emph{Journal of the American Statistical Association} \textbf{68}, 117--130 (1973).
\bibitem{breiman} Breiman, L. Stacked Regressions. \emph{Machine Learning} \textbf{24}, 49--64 (1996). \url{https://doi.org/10.1023/A:1018046112532}.
\bibitem{xu} Xu, W. H. et al. A natural language processing to causality framework for robust knowledge extraction of CO$_2$ hydrogenation with batch effect control. \emph{npj Computational Materials} \textbf{12}, 191 (2026). \url{https://doi.org/10.1038/s41524-026-02065-2}.

\bibitem{yubi} K. Yu, J. Bi and V. Tresp. Active learning via transductive experimental design. ICML, 1081--1088 (2006). \href{https://doi.org/10.1145/1143844.1143980}{doi:10.1145/1143844.1143980}.
\bibitem{fontaine} X. Fontaine, P. Perrault, M. Valko and V. Perchet. Online A-optimal design and active linear regression. ICML, PMLR 139, 3374--3383 (2021). \url{https://proceedings.mlr.press/v139/fontaine21a.html}.
\bibitem{sloman} S. J. Sloman, A. Bharti, J. Martinelli and S. Kaski. Bayesian active learning in the presence of nuisance parameters. UAI, PMLR 244, 3245--3263 (2024). \url{https://proceedings.mlr.press/v244/sloman24a.html}.
\bibitem{hooker} G. Hooker. Generalized functional ANOVA diagnostics for high-dimensional functions of dependent variables. J. Comput. Graph. Stat. 16, 709--732 (2007). \href{https://doi.org/10.1198/106186007X237892}{doi:10.1198/106186007X237892}.


\bibitem{gpml} C. E. Rasmussen and C. K. I. Williams. Gaussian Processes for Machine Learning. MIT Press (2006), Chapter 2, Section 2.7. \url{https://gaussianprocess.org/gpml/chapters/RW2.pdf}.
\bibitem{alpaca} J. Harrison, A. Sharma and M. Pavone. Meta-Learning Priors for Efficient Online Bayesian Regression. arXiv:1807.08912 (2018). \url{https://arxiv.org/abs/1807.08912}.
\bibitem{gpmllearn} Q. P. Nguyen, B. K. H. Low and P. Jaillet. Learning to learn with Gaussian processes. UAI, PMLR 161, 1466--1475 (2021). \url{https://proceedings.mlr.press/v161/nguyen21c.html}.
\bibitem{scaml} P. Tighineanu et al. Scalable Meta-Learning with Gaussian Processes. AISTATS, PMLR 238, 1981--1989 (2024). \url{https://proceedings.mlr.press/v238/tighineanu24a.html}.
\bibitem{iwata} T. Iwata and A. Kumagai. Meta-learning to Calibrate Gaussian Processes with Deep Kernels for Regression Uncertainty Estimation. arXiv:2312.07952 (2023). \url{https://arxiv.org/abs/2312.07952}.

\end{thebibliography}

\end{document}
